\section{Related Work}
\label{sec:related}

\paragraph{Prediction-powered inference.}
\citet{angelopoulos2023ppi} introduced prediction-powered inference (PPI): a small labeled sample is combined with a large unlabeled sample and the output of a fixed machine-learning predictor, and the rectifier $\bar Y_L-\bar f_L$ corrects the bias of the predictor, giving valid intervals for means, quantiles, and regression coefficients without assumptions on the predictor. \citet{angelopoulos2023ppiplus} (PPI++) add a data-driven tuning parameter $\lambda$ that interpolates between using only labels and using the predictor at full strength, and show that the resulting intervals never widen relative to labeled-only intervals in the cases they study. Our PPI++ baseline is the optimal linear calibration within this family. \citet{zrnic2023cross} cross-fit the predictor when it is trained on the same labels, and \citet{zrnic2024active} choose which points to label. Both change the sampling or fitting setting, which we hold fixed. \citet{fisch2024strat} stratify on the predictor output, and we compare to their point estimator directly (Section~\ref{sec:exp}).

\paragraph{Semiparametric efficiency with missing outcomes.}
Our efficiency bound is the standard one for estimating a mean when outcomes are missing at random with known probability, and the augmented estimator that attains it is classical \citep{robins1994estimation,tsiatis2006semiparametric}. Double/debiased machine learning \citep{chernozhukov2016dml} uses cross-fitting with a learned nuisance function to get valid inference under flexible fitting. The setting here is simpler, because the only nuisance is $m(Z)=\mathbb E[Y\mid Z]$ and it enters the estimator linearly, which lets us state the variance exactly (Proposition~\ref{prop:var}). General theory of semi-supervised mean estimation with unlabeled covariates is developed in \citet{zhang2019semi}, and semi-supervised linear regression in \citet{chakrabortty2018efficient}; we did not rely on their efficiency statements beyond the setting they describe.

\paragraph{What this paper adds.}
The question we address is narrower than the PPI literature: holding the labeled set, the unlabeled set, and the predictor fixed, can the existing intervals be made shorter without any new data or assumptions beyond validity? We answer with an explicit gap formula (Theorem~\ref{thm:bound}), a condition under which a learned calibrator helps (Corollary~\ref{cor:learned}), and a practical cross-fitted estimator whose gains we measure on real and synthetic data.
